% ============================================================================
% CHAPTER 1 — INTRODUCTION
% ============================================================================
\chapter{Introduction}
\label{chap:introduction}

\section{Clinical Context of Alzheimers disease}
\label{Clinical Context of Alzheimers disease}
Alzheimer’s disease (AD) is a progressive neurodegenerative disorder and the most common cause of dementia worldwide. Clinically, AD is characterized by a progressive decline in memory, cognitive abilities, behaviour, and the ability to perform everyday activities. Dementia represents a major global health challenge: in 2021, approximately 57 million people were living with dementia worldwide, with nearly 10 million new cases occurring each year. Alzheimer’s disease accounts for an estimated 60--70% of all dementia cases, making it one of the most important neurodegenerative diseases to investigate \cite{WHO_Dementia_2026}.

A major challenge in the management of AD is the long period between the initiation of pathological processes and the appearance of clinically recognizable symptoms. The pathological changes associated with AD can begin approximately 15--20 years before the onset of obvious cognitive symptoms \cite{Rafii_Aisen_2023}. During this preclinical phase, pathological alterations can accumulate while the affected individual may remain clinically asymptomatic. Consequently, substantial neuronal and structural changes may already be present when cognitive impairment becomes apparent. This long latent phase highlights the importance of developing sensitive methods for the early detection and characterization of AD, as early intervention may provide a greater opportunity to slow disease progression \cite{Rafii_Aisen_2023}.

In addition to the challenge of early detection, the development of effective therapies for AD remains difficult. One of the major obstacles to pharmacological treatment is the blood--brain barrier (BBB), a highly selective physiological barrier that protects the central nervous system from potentially harmful substances in the systemic circulation. While the BBB is essential for maintaining brain homeostasis, it also restricts the transport of many therapeutic compounds into the brain, thereby limiting their bioavailability at the intended site of action \cite{Iqbal_2024}. This is particularly relevant for therapeutic approaches involving large biological molecules and other compounds with limited BBB permeability.

Consequently, improving the delivery of therapeutic agents to the central nervous system has become an important research objective in AD treatment. Several strategies have been investigated to overcome or circumvent the BBB, including receptor-mediated transport, focused ultrasound, intranasal administration, liposomal systems, and nanoparticle-based drug delivery systems \cite{Iqbal_2024,DrugDelivery_BBB_AD_2025}. Nanoparticle-based approaches are particularly promising because their physicochemical properties can be engineered to improve drug solubility and stability, control drug release, enhance transport across the BBB, and potentially increase accumulation at specific target sites while reducing systemic exposure \cite{Roghani_2024,Roamcharern_2026}.

The investigation of novel drug delivery systems is therefore important for enabling therapeutic agents to reach the brain at concentrations sufficient to exert their intended effects. This is particularly relevant as increasingly sophisticated disease-modifying therapies are being developed, for which efficient and targeted delivery to the central nervous system remains a significant challenge. However, despite promising preclinical results, nanoparticle-based and other advanced delivery systems still face important translational challenges, including long-term safety, biological variability, manufacturing scalability, and regulatory requirements \cite{Roamcharern_2026}. Addressing these limitations is therefore essential for translating advanced drug delivery concepts into clinically effective therapies for Alzheimer’s disease.

Taken together, the combination of a long preclinical disease phase, the progressive and irreversible nature of neurodegeneration, and the restricted accessibility of the brain to many therapeutic compounds demonstrates the need for both improved early disease characterization and more effective drug delivery strategies. Investigating advanced drug delivery systems capable of overcoming the limitations imposed by the BBB may therefore contribute to the development of more targeted and effective therapeutic approaches for AD.


Labels of interst in the context of Alzheimer's Disease. 

\subsection{Diagnosing Alzheimers Disease}
\label{Diagnosing Alzheimers Disease}
Atrophy due to gray matter loss with anatomical magnetic resonance imaging or hypometabolism with 18F-fluorodeoxyglucose positron emission tomography (FDG PET). 

This diagnostic status is established by clinicians based on cognitive scores such as the Mini-Mental State Evaluation (MMSE),
and the Clinical Dementia Rating. Even though assessment protocols are standardized, these scores can be highly variable based on the examiner and the physical/psychological condition of the patient during the test.
\section{Foundation Models}
\label{Foundation Models}

\section{Structural MRI acquisition}
\label{sec: Structural MRI acqusition}
\textbf{Image Acqusition:}
A strong magentic fiel is applied B0, normally 1,5-3T. That forces the hydrogen nuclei of the abundant water molecuels in the soft tissue in the ody to aling with the magentic field. 
A gradient coli applies a second magentic field, which leasde to a different spinning frequency of molecules in the filed of interest. A radio-frequency (RF) impules is applied with the same 
frequency as the gradient coil made the molecules spin. This pulse which tilts these nuclei from their aligment along the B0. The nuclei then precess back to the aligment. The 
preceding nuclei emit a signal, which is registered by the reciver coils in the scanner. 
Each signal that is produced by the realignment of the hydrogen nuclei to the magnetic field B0 has two components. First the longitudinal z-axis, which is along the scanner's magentic field. 
The second is the transversal signal. This describes the xy-plane orthogonal to the scanner´s magnetic field.
\\

Initially the longitudinal signal is weak, due to the tilt of the nuclei. But over time it increases again due to the realignment of the hydrogen nuclei. The time constant 
that dictates the speed of realignment is denoted by T1. The transversal signal on the other hand is strog due to phase coherence. The signal decades as the nuclei dephase and realigne. 
This decay is denoted as the T2 time constant.\\

Tissues have different T1 and T2 relaxations times, which are used to build contrast in the MRI image. In the final image the contrast depends on the receiving time of an signal, 
the echo time (TE) and how fast you repeat the tilt-relax-process, repetition-time (TR).





\section{MRI properties for analysis}
\label{sec: MRI propperties for analysis}
Magnetic resonance Imaging (MRI) is a widely used imaging modality in radiology.
Even though MR has a long data acquisition time, and complex scanning protocols, it is preferred more than other imaging modalities such as Completed Tomography (CT)
because it lacks the use of ionizing radiation and has a better visualization potential. 

\subsection{MRI in diagnosing Alzheimer's disease}
\label{sec:MRI in diagnosing Alzheimers disease}
Alzheimer's is a complex little researched neurodegenerative disease, which can not be diagnosed within a single test. However, brain imaging plays an important role in diagnosing Alzheimer's disease.
The 3Dimage of the brain provides a great opportunity for physicians to spot abnormalities. 
These abnormalities can be markers of Alzheimer's disease, but could also be symptoms of other cognitive conditions
\cite{Envision_Radiology}. Alzheimer's is primary connected to atrophy in key regions like the hippocampus, the core memory area of the brain, and parietal lobe. 
\placeholder{
    \begin{enumerate}
        \item K-space: It is a collection of spatial frequencies captured after applying the pulse sequences.
        The (x,y) coordinates in this space represent the frequency and phase for every pixel of the image generated after applying inverse Fourier transform to the captured frequency data.
        \item Fourier transform: is a mathematical technique that allows a signal to be decomposed into a sum of sine or cosine waves of different frequencies, phases and amplitudes.
        This link can be followed for further reading about Fourier transform
    \end{enumerate}
}

\section{Foundation Models}
\label{sec:Foundation MOdels}
 Foundation models can be used as a Feature extractor 
 the workflow would be MRI-Images-> Foundation Models -> embedding matrix (Feature representation) -> Classification model -> 6 classes
 \textbf{Three ways of foundation model approches}
 \begin{enumerate}
    \item Frozen Foundation model. You dont change the weights of the model itself, instead you finetue the classifier. 
    \item Fine tune the foundation model. Fine tune is much more computational ecpensive, and can be prone to overfitting in smaller data sets
    \item Zero-shot classification. No classification needed. 
 \end{enumerate}
\section{Aim and Research Questions}
\label{sec:aim}
The aim of this project is to analyse whether pretrained foundation models can turn structural MRI slices into  useful numerical representations, then test whether those representations separate six clinical groups of alzheimers disease. 

\placeholder{State the main aim and, where appropriate, explicit research questions
or hypotheses.}
